\documentclass{beamer}

\usetheme{Madrid}

\usepackage{amsmath,amssymb,amsfonts}
\usepackage{bm}
\usepackage{physics}
\usepackage{graphicx}

\title{Stochastic Reconfiguration in Quantum Monte Carlo}
\subtitle{Variational Optimization, TDVP and Natural Gradients}
\author{Morten Hjorth-Jensen}
\date{Spring 2026}

\begin{document}

\frame{\titlepage}

%------------------------------------------------
\section{Introduction}

\begin{frame}{Quantum Many-Body Problem}

Goal: Solve the Schrödinger equation

\[
H|\Psi_0\rangle = E_0 |\Psi_0\rangle
\]

for interacting many-body systems.

Examples:

\begin{itemize}
\item Hubbard model
\item Heisenberg spin model
\item Nuclear many-body systems
\item Quantum chemistry
\end{itemize}

Exact diagonalization scales exponentially.

Need stochastic methods.
\end{frame}

%------------------------------------------------
\begin{frame}{Quantum Monte Carlo}

Quantum Monte Carlo methods include

\begin{itemize}
\item Variational Monte Carlo (VMC)
\item Diffusion Monte Carlo
\item Auxiliary-field Monte Carlo
\end{itemize}

VMC strategy:

\begin{enumerate}
\item Choose parameterized wavefunction
\item Evaluate expectation values stochastically
\item Optimize parameters
\end{enumerate}
\end{frame}

%------------------------------------------------
\begin{frame}{Variational Principle}

For any trial state

\[
|\Psi(\bm{\alpha})\rangle
\]

energy is

\[
E(\bm{\alpha}) =
\frac{\langle \Psi|H|\Psi\rangle}
{\langle \Psi|\Psi\rangle}
\]

Variational bound:

\[
E(\bm{\alpha}) \ge E_0
\]

Optimization problem:

\[
\min_{\bm{\alpha}} E(\bm{\alpha})
\]
\end{frame}

%------------------------------------------------
\section{Variational Monte Carlo}

\begin{frame}{Local Energy}

Define local energy

\[
E_L(x) =
\frac{H\Psi(x)}{\Psi(x)}
\]

Expectation value

\[
E = \frac{\int dx |\Psi(x)|^2 E_L(x)}
{\int dx |\Psi(x)|^2}
\]

Monte Carlo sampling distribution

\[
P(x) = \frac{|\Psi(x)|^2}{\langle\Psi|\Psi\rangle}
\]
\end{frame}

%------------------------------------------------
\begin{frame}{Metropolis Sampling}

Generate samples

\[
x_1,x_2,\dots,x_M
\]

using Metropolis algorithm.

Estimate expectation value

\[
E \approx
\frac{1}{M}\sum_{m=1}^{M} E_L(x_m)
\]

Statistical error

\[
\sigma_E \sim \frac{\sqrt{\text{Var}(E_L)}}{\sqrt{M}}
\]
\end{frame}

%------------------------------------------------
\section{Optimization Problem}

\begin{frame}{Energy Gradient}

Define logarithmic derivatives

\[
O_k(x)=
\frac{\partial \ln \Psi(x)}{\partial \alpha_k}
\]

Energy gradient

\[
\frac{\partial E}{\partial \alpha_k}
=
2(\langle O_k E_L\rangle
-
\langle O_k\rangle\langle E_L\rangle)
\]
\end{frame}

%------------------------------------------------
\begin{frame}{Gradient Descent}

Standard update

\[
\alpha_k^{(n+1)}
=
\alpha_k^{(n)} - \eta
\frac{\partial E}{\partial \alpha_k}
\]

Problems:

\begin{itemize}
\item Parameter correlations
\item Ill-conditioned landscape
\item Large statistical noise
\end{itemize}

Better optimization required.
\end{frame}

%------------------------------------------------
\section{Imaginary-Time Projection}

\begin{frame}{Imaginary-Time Evolution}

Ground state projection

\[
|\Psi(\tau)\rangle =
e^{-\tau H} |\Psi(0)\rangle
\]

For small step

\[
|\Psi(\tau+\delta\tau)\rangle
=
(1-\delta\tau H)|\Psi(\tau)\rangle
\]

Idea:

project evolution onto variational manifold.
\end{frame}

%------------------------------------------------
\section{Time Dependent Variational Principle}

\begin{frame}{TDVP}

TDVP minimizes distance between

\[
\partial_\tau |\Psi\rangle
\]

and

\[
-H|\Psi\rangle
\]

within variational manifold.

Resulting equation

\[
\sum_j S_{ij} \dot{\alpha}_j = -g_i
\]
\end{frame}

%------------------------------------------------
\begin{frame}{Covariance Matrix}

Define

\[
S_{ij}
=
\langle O_i O_j\rangle
-
\langle O_i\rangle\langle O_j\rangle
\]

and

\[
g_i =
\langle O_i E_L\rangle
-
\langle O_i\rangle\langle E_L\rangle
\]

Update rule

\[
\delta\alpha = -\delta\tau S^{-1} g
\]
\end{frame}

%------------------------------------------------
\section{Stochastic Reconfiguration}

\begin{frame}{SR Method}

Stochastic Reconfiguration equation

\[
\sum_j S_{ij}\delta\alpha_j = -\delta\tau g_i
\]

Interpretation:

\begin{itemize}
\item Imaginary-time evolution
\item Natural gradient descent
\item TDVP projection
\end{itemize}
\end{frame}

%------------------------------------------------
\begin{frame}{Geometric Interpretation}

Wavefunctions form a manifold.

Metric:

\[
ds^2 =
\sum_{ij} S_{ij} d\alpha_i d\alpha_j
\]

This is the

\textbf{quantum Fisher information metric}.
\end{frame}

%------------------------------------------------
\section{Quantum Information Geometry}

\begin{frame}{Quantum Fisher Information}

For normalized wavefunction

\[
|\psi(\bm{\alpha})\rangle
\]

Fisher matrix

\[
F_{ij} =
4(\langle O_i O_j\rangle
-
\langle O_i\rangle\langle O_j\rangle)
\]

Thus

\[
S_{ij} = F_{ij}/4
\]

SR corresponds to natural gradient.
\end{frame}

%------------------------------------------------
\section{Regularization}

\begin{frame}{Ill Conditioning}

Matrix \(S\) may become singular.

Regularization:

\[
S_{ij} \rightarrow S_{ij} + \lambda \delta_{ij}
\]

Typical

\[
\lambda \sim 10^{-3}
\]
\end{frame}

%------------------------------------------------
\section{Algorithm}

\begin{frame}{SR Algorithm}

\begin{enumerate}
\item Sample configurations
\item Compute \(O_k(x)\)
\item Compute local energies
\item Estimate \(S\) and \(g\)
\item Solve linear system
\item Update parameters
\end{enumerate}
\end{frame}

%------------------------------------------------
\section{Python Implementation}

\begin{frame}[fragile]{Example: Variational Monte Carlo}

\begin{verbatim}
import numpy as np

def local_energy(x, psi, Hpsi):
    return Hpsi(x)/psi(x)

def O_k(x, psi, dpsi):
    return dpsi(x)/psi(x)
\end{verbatim}

This structure is used in VMC codes.
\end{frame}

%------------------------------------------------
\begin{frame}[fragile]{SR Linear System}

\begin{verbatim}
S = np.zeros((p,p))
g = np.zeros(p)

for x in samples:

    O = compute_log_derivatives(x)

    EL = local_energy(x)

    S += np.outer(O,O)

    g += O*EL

S /= M
g /= M

delta = np.linalg.solve(S + lam*np.eye(p), -g)
\end{verbatim}

Parameter update

\begin{verbatim}
alpha += dt*delta
\end{verbatim}
\end{frame}

%------------------------------------------------
\section{Example Systems}

\begin{frame}{Heisenberg Model}

Hamiltonian

\[
H = J\sum_{\langle ij\rangle}
\mathbf{S}_i\cdot\mathbf{S}_j
\]

Typical variational states

\begin{itemize}
\item Jastrow wavefunctions
\item Neural network states
\end{itemize}
\end{frame}

%------------------------------------------------
\begin{frame}{Hubbard Model}

Hamiltonian

\[
H =
-t\sum_{\langle ij\rangle\sigma}
c^\dagger_{i\sigma}c_{j\sigma}
+
U\sum_i n_{i\uparrow}n_{i\downarrow}
\]

Applications

\begin{itemize}
\item strongly correlated electrons
\item high-Tc superconductivity
\end{itemize}
\end{frame}

%------------------------------------------------
\section{Neural Quantum States}

\begin{frame}{Neural Network Wavefunctions}

Example RBM wavefunction

\[
\Psi(\sigma)
=
\exp(\sum_i a_i\sigma_i)
\prod_j 2\cosh(b_j + W_{ij}\sigma_i)
\]

SR used to train network parameters.
\end{frame}

%------------------------------------------------
\begin{frame}{Modern Applications}

\begin{itemize}
\item Neural quantum states
\item Tensor network variational states
\item Quantum chemistry wavefunctions
\item Nuclear many-body systems
\end{itemize}
\end{frame}

%------------------------------------------------
\begin{frame}{Summary}

\begin{itemize}
\item Stochastic reconfiguration derived from TDVP
\item Equivalent to natural gradient
\item Robust optimization for VMC
\item Essential in modern neural quantum states
\end{itemize}
\end{frame}

%------------------------------------------------
\begin{frame}{References}

\begin{itemize}

\item S. Sorella, Phys. Rev. B 71 (2005)

\item F. Becca and S. Sorella  
Quantum Monte Carlo Approaches for Correlated Systems

\item G. Carleo and M. Troyer, Science (2017)

\end{itemize}

\end{frame}

\end{document}
To add
	•	full derivation of TDVP from Dirac–Frenkel principle
	•	SR vs linear method optimization
	•	complete VMC + SR code for the Heisenberg model
	•	example training of a neural quantum state
	•	connection to quantum natural gradient in quantum computing
